\documentclass[11pt]{article}
\usepackage[T1]{fontenc}
\usepackage{lmodern}
\usepackage{amsmath,amssymb,amsthm,mathtools}
\usepackage[a4paper,margin=1in]{geometry}
\usepackage{microtype}
\usepackage[hidelinks]{hyperref}
\newtheorem{theorem}{Theorem}
\newcommand{\R}{\mathbb R}
\newcommand{\F}{\mathcal F}

\title{A Certified Upper Bound for the One-Dimensional\\Fourier Sign-Uncertainty Constant}
\author{Unsolved Labs}
\date{Research Release R012, August 2026}

\begin{document}
\maketitle

\begin{abstract}
We give an explicit Fourier-invariant Schwartz function proving
\[
A_+(1)\leq\sqrt{\frac{1912071}{2000000\pi}}<0.551649
\]
for the one-dimensional $+1$ Fourier sign-uncertainty constant under the normalization
$\widehat f(y)=\int_{\R}f(x)e^{-2\pi ixy}\,dx$.
The witness is a degree-$1800$ Laguerre--Gaussian combination with $900$ frozen integer coefficients.
Its eventual positivity is certified entirely in exact integer/rational arithmetic: finite intervals are handled by Bernstein-basis positivity after exact affine changes and dyadic de Casteljau subdivision, while the infinite tail is certified by nonnegative coefficients after translation.
The search that produced the witness is outside the proof boundary.
A production C++17 verifier replays the full certificate, a separately written Python checker independently reconstructs the witness and verifies the tail and radius comparison, and a pinned Lean 4/Mathlib package kernel-checks the certificate-to-exact-radius and decimal-comparison bridge under explicit external witness obligations.
The result is an upper bound only; no matching lower bound or optimality claim is made.
\end{abstract}

\section{Introduction}
Fourier sign uncertainty asks how tightly a function and its Fourier transform can both become eventually nonnegative when they begin with the opposite sign at the origin.
In the one-dimensional $+1$ eigenfunction formulation, $A_+(1)$ measures the smallest achievable eventual-positivity radius among admissible self-Fourier Schwartz witnesses.
The rigorous one-dimensional comparison used by this release is the bound
\[
0.45\leq A_+(1)\leq0.594
\]
from Gon\c{c}alves, Oliveira e Silva, and Steinerberger (arXiv:1602.03366).
Later work, including arXiv:2003.10771, reports sharper numerical phenomena while separating those computations from rigorous proof.

R012 contributes a new explicit rigorously certified upper-bound witness.
The contribution is intentionally narrow: no claim is made that the witness is extremal, that the exact value of $A_+(1)$ is determined, or that the search procedure producing the coefficients is verified.

\section{Definitions and normalization}
We use
\[
\widehat f(y)=\int_{\R} f(x)e^{-2\pi ixy}\,dx.
\]
In the frozen $+1$ formulation, an admissible witness is a nonzero real Schwartz function with $\widehat f=f$, $f(0)\leq0$, and $f(x)\geq0$ outside a symmetric interval.
For such a witness let
\[
r(f)=\inf\{r\geq0:f(x)\geq0\text{ for every }|x|\geq r\}.
\]
Then any admissible witness gives $A_+(1)\leq r(f)$.

\section{The explicit witness}
Let $L_n^{-1/2}$ be the generalized Laguerre polynomial and define
\[
p(t)=\sum_{k=1}^{900}N_k\left(L_{2k}^{-1/2}(t)-L_{2k}^{-1/2}(0)\right),
\qquad
f(x)=p(2\pi x^2)e^{-\pi x^2},
\]
where the integers $N_1,\dots,N_{900}$ are frozen verbatim in
\texttt{coefficients/part1.txt} through \texttt{coefficients/part6.txt}.
By construction $p(0)=f(0)=0$.

The standard Laguerre--Gaussian identity is
\[
\F\!\left[L_n^{-1/2}(2\pi x^2)e^{-\pi x^2}\right]
=(-1)^nL_n^{-1/2}(2\pi x^2)e^{-\pi x^2}.
\]
Only even indices occur, hence $\widehat f=f$.
The polynomial has degree $1800$ with positive leading coefficient, so $f$ is nonzero; polynomial times Gaussian is Schwartz.
It remains to certify eventual positivity.

\section{Exact polynomial reconstruction}
Define the integer-scaled Laguerre polynomials
\[
R_n(t)=2^n n!L_n^{-1/2}(t).
\]
They satisfy
\[
R_0=1,\quad R_1=1-2t,
\]
and
\[
R_{n+1}=(4n+1-2t)R_n-2n(2n-1)R_{n-1}.
\]
The production verifier reconstructs an integer polynomial $P(t)$ that is a positive scalar multiple of $p(t)$ and checks exactly
\[
\deg P=1800,\qquad P(0)=0,
\]
with positive leading coefficient.
Set
\[
T=\frac{1912071}{1000000}=1.912071.
\]

\section{Finite-interval positivity}
For a rational interval $[a,b]$, the verifier performs an exact affine substitution to $[0,1]$ and converts the resulting polynomial to a scaled Bernstein basis.
If every Bernstein coefficient is positive, the polynomial is positive throughout the interval because the Bernstein basis functions are nonnegative on $[0,1]$.
When a parent interval is not immediately certified, exact dyadic de Casteljau subdivision is applied recursively.
All proof arithmetic is on arbitrary-precision integers; sampled positivity and floating-point root finding are absent from the proof path.

The release replay covers the complete finite interval
\[
[T,12000].
\]
The block partition in \texttt{verify\_release.sh} and CI is only computational.

\section{Infinite tail}
The verifier expands
\[
P(12000+u)=\sum_{j=0}^{1800}d_j u^j
\]
and checks exactly
\[
d_0>0,\qquad d_j\geq0\quad(1\leq j\leq1800).
\]
Therefore $P(t)>0$ for all $t\geq12000$.
Together with the finite certificates,
\[
P(t)>0\qquad(t\geq T).
\]
Since $P$ is a positive scalar multiple of $p$ and the Gaussian is positive,
\[
f(x)>0\qquad\text{whenever }2\pi x^2\geq T.
\]
Hence
\[
r(f)\leq\sqrt{\frac{T}{2\pi}}
=\sqrt{\frac{1912071}{2000000\pi}}.
\]

\section{Exact radius comparison}
The verifier does not use a floating-point value of $\pi$ as proof evidence.
It obtains a rational lower bound from Machin's identity
\[
\pi=16\arctan(1/5)-4\arctan(1/239)
\]
and alternating arctangent truncation, and proves exactly
\[
\frac{T}{2\pi}<\left(\frac{551649}{10^6}\right)^2.
\]
Thus
\[
\sqrt{\frac{T}{2\pi}}<0.551649.
\]

\begin{theorem}[R012]
Under the Fourier normalization above,
\[
A_+(1)\leq\sqrt{\frac{1912071}{2000000\pi}}<0.551649.
\]
\end{theorem}

\begin{proof}
The explicit $f$ is a nonzero Schwartz function with $f(0)=0$ and $\widehat f=f$ by the even-index Laguerre--Gaussian identity.
The exact production certificate proves $p(t)>0$ for every $t\geq T$.
Consequently $f(x)>0$ whenever $|x|\geq\sqrt{T/(2\pi)}$, so the defining variational inequality gives $A_+(1)\leq\sqrt{T/(2\pi)}$.
The exact rational $\pi$ comparison proves the final strict decimal inequality.
\end{proof}

\section{Verification architecture and trust boundary}
The repository separates discovery from verification.
The coefficient-search procedure is not required to reproduce the theorem.
A clean checkout needs only the six frozen coefficient files and the released verification artifacts.

\paragraph{Production replay.}
\texttt{verify.cpp} reconstructs the exact scaled polynomial, proves all finite Bernstein certificates, proves the translated tail, and proves the rational radius comparison.
The wrapper \texttt{verify\_release.sh} performs the complete replay; GitHub Actions shards the finite proof for operational reliability.

\paragraph{Independent partial replay.}
\texttt{verify\_independent.py} is a separately written Python implementation using arbitrary-precision integers and \texttt{fractions.Fraction}.
It independently reconstructs the degree-$1800$ witness, checks the translated infinite-tail certificate, and proves the radius comparison.
It intentionally does not duplicate the complete finite Bernstein subdivision.
Thus the finite interval certificate still has one production implementation; this limitation is public and explicit.

\paragraph{Partial Lean bridge.}
The repository pins Lean 4.32.0 and a fixed Mathlib revision.
The production declaration \texttt{R012.r012\_exact\_and\_decimal\_bounds\_from\_certificate} uses Mathlib's actual real Fourier transform to define self-Fourier admissibility.
For an abstract real function $p$, Lean proves that exact positivity for $t\geq T$, together with integrability, nontriviality, and actual Fourier self-duality of the Gaussian lift, implies
\[
A^{\mathrm{sf}}_+(1)\leq\sqrt{\frac{T}{2\pi}}<0.551649,
\]
where $A^{\mathrm{sf}}_+(1)$ is the precise self-Fourier sign-radius infimum defined in the formalization.
Lean checks the threshold-to-radius implication, exact square-root rescaling, and strict decimal comparison using a rigorous Mathlib lower bound for $\pi$.
The formal statement is isolated in \texttt{Challenge.lean}, matched against \texttt{Solution.lean} by Comparator, audited at trust level zero, replayed with \texttt{leanchecker}, and checked by the immutable reusable release contract.

This formalization is intentionally partial.
It does not kernel-check the full 900-coefficient Bernstein certificate, the generalized Laguerre--Gaussian Fourier identity for the concrete witness, the concrete integrability/nontriviality discharge, or equivalence with every alternate literature definition of $A_+(1)$.
These remaining obligations are explicit in \texttt{formalization.yaml}, \texttt{STATEMENT\_AUDIT.md}, and \texttt{manuscript/proof\_note.md}; no full Lean proof of the concrete release is claimed.

\section{Prior work and novelty boundary}
The frozen rigorous comparison is arXiv:1602.03366, supplying the $0.594$ upper bound cited by this release.
Later numerical work such as arXiv:2003.10771 is context, not a replacement rigorous theorem unless explicitly proved there.
The novelty claim is only the explicit exactly certified witness establishing the displayed $0.551649$ upper bound.
No exact value, new lower bound, or witness optimality is claimed.

\section{Reproducibility}
From a clean checkout run
\begin{verbatim}
./verify_release.sh
python3 verify_independent.py
lake exe cache get
lake build R012 Challenge Solution R012.Audit
lake env lean --trust=0 R012/Audit.lean
\end{verbatim}
The first command is the complete concrete production certificate replay; the second is the independent exact structural/tail/radius replay; the Lean commands verify the partial formal bridge.
See \texttt{CLAIM.md}, \texttt{STATEMENT\_AUDIT.md}, \texttt{VERIFICATION.md}, and \texttt{formalization.yaml} for the statement mapping and trust boundary.

\section{AI-generation disclosure}
This research artifact was generated with frontier AI and is released by Unsolved Labs with an artifact-based verification boundary.
The public claim rests on the explicit witness, exact certificate checks, the stated mathematical dependencies, and the partial kernel-checked bridge; it does not rely on private model conversations or hidden reasoning traces.

\section{References}
\begin{thebibliography}{9}
\bibitem{GOS16}
F. Gon\c{c}alves, D. Oliveira e Silva, and S. Steinerberger,
rigorous one-dimensional sign-uncertainty bounds, arXiv:1602.03366 (2016).
\bibitem{CG20}
Subsequent work on sign uncertainty and numerical structure, arXiv:2003.10771 (2020).
\end{thebibliography}
\end{document}
